% =============================================================================
% Chapter 15: Sustainable AI — ML Systems Lecture Slides (Volume II)
% =============================================================================
% IMAGES NEEDED (SVGs converted to PDF):
%   - energy-wall.pdf           - carbon-lifecycle.pdf
%   - pue-breakdown.pdf         - carbon-geography.pdf
%   - four-ms-framework.pdf     - optimization-stack.pdf
%   - jevons-paradox.pdf
% =============================================================================
\documentclass[aspectratio=169, 12pt]{beamer}
\usepackage{../../assets/beamerthememlsys}

\mlsyssetup{
  volume       = {Volume II},
  chapter      = {Chapter 15},
  logo         = {../../assets/img/logo-mlsysbook.png},
  instlogo     = {../../assets/img/logo-harvard.png},
  chaptertitle = {Sustainable AI},
}

% --- Fonts ---
\usepackage{fontspec}
\setsansfont{Helvetica Neue}[
  BoldFont={Helvetica Neue Bold},
  ItalicFont={Helvetica Neue Italic},
  BoldItalicFont={Helvetica Neue Bold Italic},
]
\setmonofont{JetBrains Mono}[Scale=0.85]

% --- Packages ---
\usepackage{booktabs}
\usepackage{amsmath}

% --- Image paths ---
\graphicspath{
  {images/}
}

% --- Chapter-specific macros ---
\newcommand{\COtwo}{CO\textsubscript{2}}

% --- Helper: safe image include ---
\newcommand{\safeimg}[2][width=\textwidth,keepaspectratio]{%
  \IfFileExists{images/#2}{\includegraphics[#1]{#2}}{%
    \IfFileExists{#2}{\includegraphics[#1]{#2}}{%
      \fbox{\parbox[c][2.5cm][c]{0.85\linewidth}{\centering\footnotesize\textcolor{midgray}{[Missing image]}}}%
    }%
  }%
}

% --- Section count for navigation (must match actual \section{} count) ---
\setcounter{mlsystotalsections}{7}

\title{Sustainable AI}
\author{Vijay Janapa Reddi}
\institute{Harvard University}
\date{}

\begin{document}

% =============================================================================
% TITLE SLIDE
% =============================================================================
\mlsystitle{Sustainable AI}{Energy as a First-Class Engineering Constraint}{cover_sustainable_ai.png}

% =============================================================================
% LEARNING OBJECTIVES
% =============================================================================
\begin{frame}{Learning Objectives}
\note{
% -- NARRATE: What to SAY while showing this slide
[2 min] Walk through each objective. ``This chapter treats energy as a hard
physical constraint --- like memory or bandwidth --- not a corporate
responsibility metric.''

% -- ENGAGE: Specific question, prediction, or task for THIS slide
Ask: ``How many of you have ever checked the carbon footprint of a training run?''
[Expected: very few hands. This motivates the chapter.]

% -- FLEX: [CORE] + contingency
[CORE] Sets the chapter scope. IF SHORT: Read objectives 1, 2, and 5 aloud.
}

\small
\begin{enumerate}
  \item Explain the \textbf{sustainability paradox}: AI compute growth (350,000$\times$) outpaces hardware efficiency gains
  \item Calculate \textbf{PUE} and \textbf{lifecycle carbon footprints} across training, inference, and manufacturing phases
  \item Analyze geographic and temporal factors affecting \textbf{carbon intensity} across energy grids
  \item Evaluate \textbf{pruning, quantization, and distillation} in terms of energy trade-offs
  \item Design \textbf{carbon-aware scheduling} strategies for 50--80\% emission reductions
  \item Critique \textbf{carbon offsets} versus actual emissions reduction strategies
\end{enumerate}

\end{frame}

\begin{frame}{Visual Language}
\note{
% -- NARRATE: What to SAY while showing this slide
[1 min] Point to each card. ``In this chapter, red marks the energy cost and
carbon bottleneck. Green marks efficiency gains and clean energy paths.''

% -- FLEX: [OPTIONAL] + contingency
[OPTIONAL] Skip if students have seen this in a previous lecture.
}

\small
Throughout this course, colors carry meaning:

\vspace{0.3cm}
\begin{columns}[T]
  \begin{column}{0.45\textwidth}
    \begin{mlsyscard}{computestroke}
    \textbf{Blue} --- Compute / Processing\\
    {\footnotesize GPU ops, forward/backward pass, inference}
    \end{mlsyscard}
    \vspace{0.15cm}
    \begin{mlsyscard}{datastroke}
    \textbf{Green} --- Data / Memory\\
    {\footnotesize Data flow, caches, healthy paths}
    \end{mlsyscard}
  \end{column}
  \begin{column}{0.45\textwidth}
    \begin{mlsyscard}{routingstroke}
    \textbf{Orange} --- Routing / Scheduling\\
    {\footnotesize Load balancers, batch windows}
    \end{mlsyscard}
    \vspace{0.15cm}
    \begin{mlsyscard}{errorstroke}
    \textbf{Red} --- Error / Cost / Bottleneck\\
    {\footnotesize Loss, decode phase, waste}
    \end{mlsyscard}
  \end{column}
\end{columns}
\end{frame}



% =============================================================================
\section{The Energy Ceiling}
% =============================================================================

\begin{frame}{The Scale of the Problem}
\note{
% -- LINK: What prior concept connects to this slide
Ch 14 showed robustness costs compute. This chapter asks: what does all that
compute cost the planet?

% -- NARRATE: What to SAY while showing this slide
[3 min] Point to the left column: ``Training GPT-3: 1,287 MWh, 552 metric
tonnes CO2 --- equivalent to 552 transatlantic flights.'' Point to the table
on the right: ``Transformer+NAS: 284,000 kg --- 57 times a US car's lifetime
carbon.'' Key reveal: ``AI compute grew 350,000x from 2012--2019. Hardware
efficiency grew only 2--5\% per year. The gap is physical.''

% -- ENGAGE: Specific question, prediction, or task for THIS slide
Ask: ``What other industry produces this much carbon from a single product?''
[Expected: steel, cement, aviation. Point out datacenters are approaching all three.]

% -- WARN: What students will get wrong on THIS topic
Common error: students assume cloud = efficient = sustainable. Correct framing:
cloud is more efficient than on-premises but still consumes enormous energy.

% -- FLEX: [CORE] + contingency
[CORE] The 350,000x gap is the chapter's motivating number.
IF SHORT: State the GPT-3 number and the 350,000x gap.
}

\small
\begin{columns}[T]
  \begin{column}{0.55\textwidth}
    Training a single frontier model:
    \begin{itemize}\setlength\itemsep{0pt}
      \item \textbf{Energy}: 1,287 MWh (GPT-3)
      \item \textbf{Equivalent}: 122 US homes for one year
      \item \textbf{Carbon}: 552 metric tons \COtwo{}
      \item \textbf{Comparison}: 552 transatlantic flights
    \end{itemize}

    \vspace{0.15cm}
    \begin{mlsyscard}{errorstroke}
    AI compute grew \textbf{350,000$\times$} from 2012--2019.\\
    Hardware efficiency grew only \textbf{2--5\%/year}.
    \end{mlsyscard}
  \end{column}
  \begin{column}{0.42\textwidth}
    \scriptsize
    \renewcommand{\arraystretch}{1.1}
    \begin{tabular}{@{}lr@{}}
      \toprule
      \textbf{Activity} & \textbf{kg \COtwo{}} \\
      \midrule
      NY--SF flight      & 900 \\
      Human (1 year avg) & 5,000 \\
      American (1 year)  & 16,400 \\
      US car (lifetime)  & 57,150 \\
      \textcolor{errorstroke}{\textbf{Transformer+NAS}} & \textcolor{errorstroke}{\textbf{284,000}} \\
      \bottomrule
    \end{tabular}
  \end{column}
\end{columns}
\mlsyscite{Patterson et al., 2021; Strubell et al., 2019}

\end{frame}

\begin{frame}{The Energy Wall}
\note{
% -- LINK: What prior concept connects to this slide
The scale numbers showed the gap. This diagram visualizes WHY the gap exists ---
exponential demand vs.\ linear supply.

% -- NARRATE: What to SAY while showing this slide
[3 min] Point to the two curves: ``AI compute scales exponentially. Energy
infrastructure scales linearly. The 195,000x gap cannot be closed by hardware
alone.'' This is the Energy Wall --- physics, not budget, is the binding constraint.

% -- ENGAGE: Specific question, prediction, or task for THIS slide
Ask: ``Which curve can we change?'' [Expected: the demand curve via efficiency.
Also accept: the supply curve via renewables. Both are needed.]

% -- FLEX: [CORE] + contingency
[CORE] The Energy Wall diagram is the chapter's anchor visual.
IF SHORT: State ``exponential vs.\ linear'' and the 195,000x gap.
}

% --- Layout: FULL-WIDTH IMAGE ---
\centering
\safeimg[width=0.92\textwidth,height=4.5cm,keepaspectratio]{energy-wall.pdf}

\vspace{0.15cm}
\mlsysalert{The Sustainability Wall:}{AI logic scales exponentially; energy infrastructure scales linearly. The \textbf{195,000$\times$ gap} cannot be closed by hardware alone.}

\end{frame}

\begin{frame}{Datacenter Energy Projections}
\note{
% -- LINK: What prior concept connects to this slide
The Energy Wall showed the physics. This slide shows the infrastructure reality:
grid capacity is the binding constraint.

% -- NARRATE: What to SAY while showing this slide
[2 min] Walk through the projections: ``US datacenters: 58 TWh in 2014, 176 TWh
in 2023, projected 580 TWh by 2028 --- 12\% of US electricity.'' Point to the
table: ``Datacenters are approaching aviation and cement in global power share.''

% -- WARN: What students will get wrong on THIS topic
Common error: students think this is a cost problem. Correct framing: it is a
grid capacity problem --- you cannot buy electricity that does not exist.

% -- FLEX: [OPTIONAL] + contingency
[OPTIONAL] Extends the Energy Wall with concrete projections.
IF SHORT: State the 12\% US electricity projection and move on.
}

\small
\begin{columns}[T]
  \begin{column}{0.55\textwidth}
    \textbf{US Datacenter Electricity (LBNL 2024):}
    \begin{itemize}\setlength\itemsep{0pt}
      \item 2014: 58 TWh (1.7\% of US)
      \item 2023: 176 TWh (4.4\% of US)
      \item 2028 (high): 580 TWh (\textbf{12\% of US})
    \end{itemize}

    \vspace{0.15cm}
    AI workloads: 10\% $\to$ 30\%+ of DC power by 2030.

    \vspace{0.1cm}
    \begin{mlsyscard}{crimson}
    \textbf{Grid capacity} --- not budget --- is the binding constraint on AI scaling.
    \end{mlsyscard}
  \end{column}
  \begin{column}{0.42\textwidth}
    \scriptsize
    \renewcommand{\arraystretch}{1.1}
    \begin{tabular}{@{}lr@{}}
      \toprule
      \textbf{Sector} & \textbf{\% Global Power} \\
      \midrule
      Aviation         & 2.1\% \\
      Data centers     & 2\% (today) \\
      Data centers     & \textbf{8\%} (2030) \\
      Cement           & 4\% \\
      \bottomrule
    \end{tabular}
  \end{column}
\end{columns}
\mlsyscite{LBNL 2024; OECD 2023}

\end{frame}

% --- ACTIVE LEARNING 1: Predict ---
\begin{frame}{Predict: Where Does the Carbon Hide?}
\note{
% -- LINK: What prior concept connects to this slide
Students saw the scale of training carbon. This prediction tests whether they
recognize that inference often dominates at production scale.

% -- NARRATE: What to SAY while showing this slide
[2 min] Read the scenario: ``A model trained once, serving 10M queries/day.
Which phase produces more total CO2 over one year?'' Give 60 seconds. Do NOT
reveal yet. Ask 2--3 students to share. Most will say training.

% -- ENGAGE: Specific question, prediction, or task for THIS slide
Falsifiable prediction: training vs.\ inference carbon. [Expected: most say
training. The surprise reveal on the next slide: inference exceeds training
within days of deployment at this query volume.]

% -- FLEX: [CORE] + contingency
[CORE] First active learning moment. Primes the lifecycle analysis.
IF SHORT: Reduce writing to 30 seconds; still cold-call 2 students.
}

\centering
\vspace{1.0cm}
{\Large\bfseries Think--Write--Share}

\vspace{0.5cm}
{\large A model is trained once, then serves 10M queries/day.\\[0.2cm]
\alert{Which phase produces more total \COtwo{} over one year?}}

\vspace{0.5cm}
{\normalsize Write your answer with a rough estimate. \textcolor{midgray}{(60 seconds)}}

\vspace{0.4cm}
{\small\textcolor{lightgray}{Turn to a neighbor and compare.}}

\end{frame}

% =============================================================================
\section{Carbon Accounting}
% =============================================================================

\begin{frame}{Carbon Lifecycle: Three Phases}
\note{
% -- LINK: What prior concept connects to this slide
Most students predicted training dominates. This slide reveals the surprise:
inference overtakes training within days at production scale.

% -- NARRATE: What to SAY while showing this slide
[3 min] ``How many of you said training? Most people do.'' Walk through the
diagram: ``Training: 60--80\% of initial carbon. Inference: 15--25\%.
Manufacturing: 5--15\%. But at 10M queries/day, inference exceeds training
within 2 days of deployment.''

% -- WARN: What students will get wrong on THIS topic
Common error: students optimize training carbon and ignore inference. Correct
framing: at production scale, inference is the dominant sustainability lever.

% -- FLEX: [CORE] + contingency
[CORE] The ``2 days'' reveal reframes the entire optimization target.
}

% --- Layout: FULL-WIDTH IMAGE ---
\centering
\safeimg[width=0.92\textwidth,height=4.2cm,keepaspectratio]{carbon-lifecycle.pdf}

\vspace{0.1cm}
\mlsysinsight{Surprise:}{At production scale, inference emissions exceed training within \textbf{2 days} of deployment.}

\end{frame}

\begin{frame}{Operational Carbon Calculation}
\note{
% -- LINK: What prior concept connects to this slide
The lifecycle showed three phases. This worked example calculates operational
carbon for a real training run, revealing geographic leverage.

% -- NARRATE: What to SAY while showing this slide
[3 min] Walk through the table row by row using dimensional analysis: ``64 GPUs
at 400W for 336 hours = 8,602 kWh. Multiply by PUE 1.2 = 10,322 kWh. On the
US grid at 429 g/kWh = 4,428 kg CO2. On Quebec hydro at 20 g/kWh = 206 kg.''
Key reveal: ``Same model, same hardware, same time --- 21x difference from
geography alone.''

% -- WARN: What students will get wrong on THIS topic
Common error: students forget the PUE multiplier. Correct framing: facility
overhead (cooling, networking) adds 20--58\% on top of IT power.

% -- FLEX: [CORE] + contingency
[CORE] The 21x geographic multiplier is the chapter's key quantitative anchor.
IF SHORT: Show the final two rows (US vs.\ Quebec) and state the 21x.
}

\small
\textbf{Training a 7B model: 64 A100 GPUs, 14 days}

\vspace{0.1cm}
\renewcommand{\arraystretch}{1.15}
{\footnotesize
\begin{tabular}{@{}llll@{}}
  \toprule
  \textbf{Step} & \textbf{Calculation} & \textbf{Result} \\
  \midrule
  \textcolor{computestroke}{GPU Energy}    & 64 $\times$ 400W $\times$ 336h  & 8,602 kWh \\
  \textcolor{routingstroke}{Facility (PUE 1.2)} & 8,602 $\times$ 1.2        & 10,322 kWh \\
  \textcolor{errorstroke}{US Grid (429 g/kWh)} & 10,322 $\times$ 0.429    & \textbf{4,428 kg \COtwo{}} \\
  \textcolor{datastroke}{Quebec (20 g/kWh)}    & 10,322 $\times$ 0.020    & \textbf{206 kg \COtwo{}} \\
  \bottomrule
\end{tabular}
}

\vspace{0.15cm}
\begin{mlsyscard}{crimson}
\textbf{Geographic choice alone}: 21$\times$ difference in emissions.\\
{\footnotesize Same model, same hardware, same training time --- only the grid changed.}
\end{mlsyscard}

\end{frame}

\mlsysfocus{The Carbon Equation}{%
$C_{lifecycle} = \underbrace{C_{training}}_{\text{60--80\%}} +
\underbrace{C_{inference}}_{\text{15--25\%}} +
\underbrace{C_{embodied}}_{\text{5--15\%}}$%
}

\begin{frame}{The Geography of Carbon}
\note{
% -- LINK: What prior concept connects to this slide
The worked example showed 21x from US vs.\ Quebec. This map reveals the full
40x range across global grids.

% -- NARRATE: What to SAY while showing this slide
[3 min] Walk through the map: ``Quebec hydro at 20 g/kWh vs.\ Poland coal at
820 g/kWh = 40x difference. This is larger than ANY algorithmic optimization.''
Key statement: ``Site selection is the single highest-leverage sustainability
intervention.''

% -- ENGAGE: Specific question, prediction, or task for THIS slide
Ask: ``If you could only change one thing, would you optimize the model or move
it?'' [Expected: move it --- 40x vs.\ 2--10x from algorithmic optimization.]

% -- FLEX: [CORE] + contingency
[CORE] The 40x geographic leverage is the chapter's most actionable insight.
}

% --- Layout: FULL-WIDTH IMAGE ---
\centering
\safeimg[width=0.90\textwidth,height=4.5cm,keepaspectratio]{carbon-geography.pdf}

\vspace{0.1cm}
\small
\textbf{Site selection is the single highest-leverage sustainability intervention.}

\end{frame}

% --- ACTIVE LEARNING: Micro-Retrieval Cue ---
\begin{frame}{Quick Check}
\note{
% -- NARRATE: What to SAY while showing this slide
[1 min] Read the question. Give 15 seconds. Cold-call. Answer: geography wins ---
40x vs.\ 2--10x from algorithmic optimization.

% -- FLEX: [CORE] + contingency
[CORE] Micro-retrieval reinforcing the key geographic leverage number.
}

\centering
\vspace{1.0cm}
{\Large\bfseries Quick Check}

\vspace{0.6cm}
{\large What creates a larger carbon reduction:\\optimizing the model or moving it to a clean grid?}

\vspace{0.5cm}
{\normalsize\textcolor{midgray}{15 seconds --- then I will cold-call.}}

\end{frame}



% =============================================================================
\section{Power Modeling}
% =============================================================================

\begin{frame}{CMOS Power Equation}
\note{
% -- LINK: What prior concept connects to this slide
Geographic optimization addresses WHERE. Power modeling addresses HOW ---
grounding efficiency in CMOS physics.

% -- NARRATE: What to SAY while showing this slide
[3 min] Write the equation on screen. ``The V-squared term is the key lever:
reducing voltage from 1.0V to 0.8V saves 36\% dynamic power. This explains
why specialized accelerators (high alpha, low V) beat CPUs (low alpha, high V)
on energy efficiency.''

% -- ENGAGE: Specific question, prediction, or task for THIS slide
Ask: ``Why does reducing voltage by 20\% save 36\% power, not 20\%?''
[Expected: because power scales as V-squared --- quadratic, not linear.]

% -- WARN: What students will get wrong on THIS topic
Common error: students assume power scales linearly with voltage. Correct
framing: the V-squared dependence means small voltage reductions yield
disproportionate power savings.

% -- FLEX: [CORE] + contingency
[CORE] The CMOS equation explains WHY quantization and specialization save energy.
IF SHORT: State the V-squared insight and the 36\% example.
}

\small
$$P_{total} = \underbrace{\alpha C V^2 f}_{\text{Dynamic}} + \underbrace{V \cdot I_{leak}}_{\text{Static}}$$

\vspace{0.05cm}
\renewcommand{\arraystretch}{1.05}
{\footnotesize
\begin{tabular}{@{}ll@{}}
  \toprule
  \textbf{Param} & \textbf{Meaning} \\
  \midrule
  $\alpha$ & Switching activity (CPU: 0.1--0.3, Accel: 0.6--0.8) \\
  $C$      & Load capacitance (scales with transistor count) \\
  $V$      & Supply voltage (\textbf{quadratic} --- highest lever) \\
  $f$      & Clock frequency (linear) \\
  \bottomrule
\end{tabular}
}

\vspace{0.05cm}
\mlsysconcept{Key:}{$V^2$ means 20\% voltage reduction saves 36\% power. Specialization wins.}

\end{frame}

\begin{frame}{PUE: The Cost of Cooling}
\note{
% -- LINK: What prior concept connects to this slide
CMOS physics set the chip-level floor. PUE shows how much the facility adds
on top --- the cooling tax.

% -- NARRATE: What to SAY while showing this slide
[2 min] Point to the equation: ``PUE = total facility power / IT equipment
power. Industry average 1.58 means 37\% goes to cooling.'' Point to the table:
``Google best: 1.08 --- only 8\% overhead. For a 2 MW cluster, the gap between
1.58 and 1.08 = 8,400 MWh and \$589K saved per year.''

% -- FLEX: [OPTIONAL] + contingency
[OPTIONAL] Extends the CMOS slide to facility-level.
IF SHORT: State the PUE formula and the 1.58 vs.\ 1.08 comparison.
}

% --- Layout: FULL-WIDTH IMAGE + compact summary ---
\centering
\safeimg[width=0.85\textwidth,height=3.2cm,keepaspectratio]{pue-breakdown.pdf}

\vspace{0.05cm}
\scriptsize
$PUE = P_{total\_facility}\;/\;P_{IT\_equipment}$ \quad (ideal = 1.0)
\renewcommand{\arraystretch}{0.9}
\begin{tabular}{@{}lrrr@{}}
  \toprule
  \textbf{Facility} & \textbf{PUE} & \textbf{Overhead} & \textbf{2 MW Savings} \\
  \midrule
  Industry average & 1.58 & 37\% & --- \\
  Google best      & 1.08 & 8\%  & 8,409 MWh/yr \\
  \bottomrule
\end{tabular}

\end{frame}

\begin{frame}{Energy Per Operation: 4 Orders of Magnitude}
\note{
% -- LINK: What prior concept connects to this slide
CMOS physics explained WHY specialization saves power. This table quantifies
HOW MUCH across the architecture spectrum.

% -- NARRATE: What to SAY while showing this slide
[2 min] Walk through the table: ``CPU: 100 pJ/FLOP. GPU tensor cores: 10.
TPU systolic array: 1--2. Fixed-function ASIC: 0.1. That is a 1,000x range.''
Key insight: ``Choosing the right hardware for the workload is free
sustainability --- no algorithmic changes needed.''

% -- FLEX: [OPTIONAL] + contingency
[OPTIONAL] Reinforces the CMOS physics with concrete numbers.
IF SHORT: State the 1,000x range and the green card insight.
}

\small
\renewcommand{\arraystretch}{1.15}
{\footnotesize
\begin{tabular}{@{}lrl@{}}
  \toprule
  \textbf{Architecture} & \textbf{pJ/FLOP} & \textbf{Characteristics} \\
  \midrule
  \textbf{CPU (general)}     & 100   & Low utilization, high flexibility \\
  \textbf{GPU (tensor cores)} & 10   & High throughput, parallel \\
  \textbf{TPU (systolic)}     & 1--2 & Specialized dataflow \\
  \textbf{ASIC (INT8)}        & 0.1  & Fixed-function, low precision \\
  \bottomrule
\end{tabular}
}

\vspace{0.2cm}
\begin{mlsyscard}{datastroke}
\textbf{1,000$\times$ efficiency range} across architectures. Choosing the right hardware
tier for a workload can reduce energy consumption by 100--1,000$\times$ without any
algorithmic changes.
\end{mlsyscard}

\end{frame}

\begin{frame}{Energy Roofline: Arithmetic Intensity Meets Energy}
\note{
% -- LINK: What prior concept connects to this slide
The energy-per-operation table showed 1,000x across architectures. The energy
roofline model explains WHY --- it maps workload arithmetic intensity to
energy efficiency, just as the compute roofline maps to throughput.

% -- NARRATE: What to SAY while showing this slide
[3 min] Point to the diagram: ``X-axis is arithmetic intensity (FLOPs/byte).
Y-axis is energy efficiency (FLOPs/Joule). Below the bandwidth roof: memory-bound,
energy wasted moving data. Above: compute-bound, energy goes to useful work.''
Walk through the three workloads: ``Embedding lookup: 0.1 FLOPs/byte --- deeply
memory-bound, wasting 90\% of energy on data movement. Attention (batch=1):
2 FLOPs/byte --- straddling the roofline. GEMM (batch=256): 100 FLOPs/byte ---
compute-bound, energy-efficient.''

% -- ENGAGE: Specific question, prediction, or task for THIS slide
Ask: ``Which optimization --- batching, quantization, or pruning --- moves a
workload rightward on this plot?'' [Expected: batching increases arithmetic
intensity by amortizing memory access across more compute.]

% -- WARN: What students will get wrong on THIS topic
Common error: students think energy efficiency is purely a hardware property.
Correct framing: the same hardware can be 10--50x more energy-efficient
depending on the workload's arithmetic intensity --- software determines
where you sit on the roofline.

% -- FLEX: [CORE] + contingency
[CORE] The energy roofline connects all optimization techniques to physics.
IF SHORT: Name the three workloads and their positions on the roofline.
}

\small
\begin{columns}[T]
  \begin{column}{0.55\textwidth}
    \textbf{Energy Roofline Model:}

    \vspace{0.1cm}
    {\footnotesize Like the compute roofline, but Y-axis is \textbf{FLOPs/Joule} instead of FLOPs/sec.}

    \vspace{0.1cm}
    \scriptsize
    \renewcommand{\arraystretch}{1.1}
    \begin{tabular}{@{}lrl@{}}
      \toprule
      \textbf{Workload} & \textbf{FLOPs/B} & \textbf{Regime} \\
      \midrule
      Embedding lookup & 0.1  & \textcolor{errorstroke}{\textbf{Mem-bound}} \\
      Attention (B=1)  & 2    & \textcolor{routingstroke}{\textbf{Transitional}} \\
      GEMM (B=256)     & 100  & \textcolor{datastroke}{\textbf{Compute-bound}} \\
      \bottomrule
    \end{tabular}

    \vspace{0.1cm}
    \mlsysconcept{Lever:}{Move workloads \textbf{rightward} (higher AI) to improve energy efficiency.}
  \end{column}
  \begin{column}{0.42\textwidth}
    \begin{mlsyscard}{errorstroke}
    \textbf{Memory-bound = Energy waste}\\[0.05cm]
    {\scriptsize 90\%+ of energy moves data, not computes. LLM decode is the worst offender: reads full weights per token.}
    \end{mlsyscard}

    \vspace{0.1cm}
    \begin{mlsyscard}{datastroke}
    \textbf{Optimization strategies:}\\[0.05cm]
    {\scriptsize
    Batching: $\uparrow$ AI by $B\times$\\
    Quantization: $\uparrow$ AI by 2--4$\times$\\
    FlashAttention: $\uparrow$ AI by fusing ops\\
    All shift workloads toward the compute roof.}
    \end{mlsyscard}
  \end{column}
\end{columns}

\end{frame}

% --- ACTIVE LEARNING 2: Exercise ---
\begin{frame}{Your Turn: Lifecycle Carbon Estimation}
\note{
% -- LINK: What prior concept connects to this slide
The worked example calculated operational carbon. This exercise adds embodied
carbon to compute the full lifecycle footprint.

% -- NARRATE: What to SAY while showing this slide
[3 min] Read the parameters. Give 90 seconds. After pause reveal: ``Operational:
2048 x 0.7 x 1.3 x 720h = 1.34M kWh x 400g = 536 tonnes. Embodied:
2048 x 1500 / 36 months = 85 tonnes. Total = 621 tonnes.''

% -- ENGAGE: Specific question, prediction, or task for THIS slide
Students must calculate both operational and embodied carbon. Key insight to
surface: embodied is 14\% --- significant, and dominant on clean grids.

% -- FLEX: [CORE] + contingency
[CORE] Second active learning moment. IF SHORT: Reduce to 60 seconds.
}

\small
\begin{columns}[T]
  \begin{column}{0.56\textwidth}
    {\normalsize\bfseries Calculate Total Carbon Footprint}

    \vspace{0.15cm}
    \textbf{Parameters:}
    \begin{itemize}\setlength\itemsep{0pt}
      \item 2,048 H100 GPUs, 30 days, 700 W TDP
      \item PUE = 1.3, Grid = 400 g \COtwo{}/kWh
      \item Embodied: 1,500 kg \COtwo{}/chip
      \item Hardware lifetime: 3 years
    \end{itemize}

    \vspace{0.1cm}
    \textbf{Calculate}: Operational + Embodied = Total

    {\footnotesize\textcolor{midgray}{(90 seconds --- then compare with a neighbor)}}
  \end{column}
  \begin{column}{0.40\textwidth}
    \pause
    \begin{mlsyscard}{computestroke}
    \textbf{Solution:}\\[0.1cm]
    {\footnotesize
    Oper: $2048 \times 0.7 \times 1.3 \times 720$h\\
    = 1.34M kWh $\times$ 400g\\
    = \textbf{536 tonnes}\\[0.1cm]
    Embodied: $\frac{2048 \times 1500}{36}$\\
    = \textbf{85 tonnes}\\[0.1cm]
    Total = \textbf{621 tonnes \COtwo{}}\\[0.1cm]
    \textcolor{routingstroke}{Embodied = 14\%}
    }
    \end{mlsyscard}
  \end{column}
\end{columns}

\end{frame}

% =============================================================================
\section{Optimization Techniques}
% =============================================================================

\begin{frame}{The Optimization Toolkit}
\note{
% -- LINK: What prior concept connects to this slide
The CMOS equation showed P = alpha C V-squared f. This slide connects three
optimization techniques to the specific terms they reduce.

% -- NARRATE: What to SAY while showing this slide
[3 min] Point to the diagram: ``Pruning reduces C by removing weights.
Quantization reduces C (narrower datapaths) and V (larger noise margins).
Distillation reduces the entire model.'' Key insight: ``These are multiplicative
when combined --- pruning x quantization x distillation.''

% -- FLEX: [CORE] + contingency
[CORE] Maps optimizations to physics. IF SHORT: Name the three techniques and
their CMOS terms.
}

% --- Layout: FULL-WIDTH IMAGE ---
\centering
\safeimg[width=0.92\textwidth,height=4.5cm,keepaspectratio]{optimization-stack.pdf}

\vspace{0.15cm}
\small
\textbf{All techniques reduce terms in:} $P = \alpha C V^2 f + V \cdot I_{leak}$

\end{frame}

\begin{frame}{Quantization: The Highest-Leverage Technique}
\note{
% -- LINK: What prior concept connects to this slide
The optimization stack showed quantization reduces C and V. This slide quantifies
the savings across precision levels.

% -- NARRATE: What to SAY while showing this slide
[2 min] Walk through the table: ``FP32 baseline. FP16: 4x energy, 2x memory.
INT8: 16x energy, 4x memory. INT4: 64x energy, 8x memory.'' Point to the card:
``Narrower datapaths reduce C by 4x, larger noise margins reduce V by 2x,
simpler control reduces logic by 2x. Combined: 6--10x per operation.''
Concrete example: ``GPTQ enables 4-bit LLaMA-65B: 130 GB to 32 GB ---
consumer GPU deployment.''

% -- FLEX: [CORE] + contingency
[CORE] Quantization is the highest-leverage single technique.
IF SHORT: State the INT8 = 16x number and the GPTQ example.
}

\footnotesize
\begin{columns}[T]
  \begin{column}{0.50\textwidth}
    \textbf{Energy savings by precision:}

    \vspace{0.1cm}
    \renewcommand{\arraystretch}{1.15}
    \begin{tabular}{@{}lrr@{}}
      \toprule
      \textbf{Precision} & \textbf{Energy/Op} & \textbf{Memory} \\
      \midrule
      FP32 (baseline) & 1$\times$  & 1$\times$ \\
      FP16 / BF16     & 4$\times$  & 2$\times$ \\
      INT8            & 16$\times$ & 4$\times$ \\
      INT4            & 64$\times$ & 8$\times$ \\
      \bottomrule
    \end{tabular}
  \end{column}
  \begin{column}{0.46\textwidth}
    \begin{mlsyscard}{datastroke}
    \textbf{Why it works:}\\[0.05cm]
    {\scriptsize
    Narrower datapaths reduce $C$ (4$\times$)\\
    Larger noise margins reduce $V$ (2$\times$)\\
    Simpler control reduces logic (2$\times$)\\[0.1cm]
    Combined: \textbf{6--10$\times$ per operation}}
    \end{mlsyscard}

    \vspace{0.1cm}
    \textbf{GPTQ}: 4-bit LLaMA-65B\\
    130 GB $\to$ 32 GB\\
    \textcolor{datastroke}{\textbf{Consumer GPU deployment}}
  \end{column}
\end{columns}

\end{frame}

\begin{frame}{Training vs.\ Inference Energy}
\note{
% -- LINK: What prior concept connects to this slide
Quantization optimizes per-operation energy. This slide shows where that
optimization matters most: inference dominates at production scale.

% -- NARRATE: What to SAY while showing this slide
[2 min] Point to the left cards: ``Training: one-time burst, 1,287 MWh for
GPT-3. Inference: continuous drain, 10M queries/day = 3,650 MWh/year.''
ANALOGY: ``Training is the rocket launch; inference is the airline fleet.''
Point to the decode problem: ``Decode reads full weights per token --- 10--50x
less energy-efficient than batch processing.''

% -- WARN: What students will get wrong on THIS topic
Common error: students optimize training energy because it is visible and
one-time. Correct framing: inference exceeds training within 2 days at 10M
queries/day --- inference optimization is the dominant lever.

% -- FLEX: [CORE] + contingency
[CORE] Reframes the optimization target from training to inference.
IF SHORT: State the 2-day crossover and the decode inefficiency.
}

\small
\begin{columns}[T]
  \begin{column}{0.48\textwidth}
    \begin{mlsyscard}{errorstroke}
    \textbf{Training}: One-time burst\\[0.05cm]
    {\footnotesize 1,287 MWh for GPT-3\\
    4.4 tonnes \COtwo{} (US grid)}
    \end{mlsyscard}

    \vspace{0.15cm}
    \begin{mlsyscard}{computestroke}
    \textbf{Inference}: Continuous drain\\[0.05cm]
    {\footnotesize 10M queries/day $\times$ 0.001 kWh\\
    = 10,000 kWh/day = 3,650 MWh/year\\
    = 1,566 tonnes \COtwo{}/year}
    \end{mlsyscard}
  \end{column}
  \begin{column}{0.48\textwidth}
    \vspace{0.3cm}
    \textbf{The Decode Problem:}

    \vspace{0.1cm}
    {\footnotesize
    \textbf{Prefill} (compute-bound):\\
    High arithmetic intensity $\to$ efficient\\[0.1cm]
    \textbf{Decode} (bandwidth-bound):\\
    Reads full weights per token $\to$ GPU idles\\[0.1cm]
    \textcolor{errorstroke}{\textbf{10--50$\times$ less energy-efficient}}\\
    than batch processing
    }

    \vspace{0.1cm}
    \mlsysconcept{Lever:}{Optimize inference, not just training.}
  \end{column}
\end{columns}

\end{frame}

% =============================================================================
\section{Carbon-Aware Scheduling}
% =============================================================================

\begin{frame}{Carbon-Aware Scheduling}
\note{
% -- LINK: What prior concept connects to this slide
Geography showed 40x across regions. Carbon-aware scheduling exploits that
variation dynamically --- shifting jobs in time and space.

% -- NARRATE: What to SAY while showing this slide
[3 min] Point to the table: ``Coal: 820--1,200 g/kWh. Hydro: 10--30 g/kWh.
Grid intensity varies 100x.'' Point to the two strategies: ``Temporal: delay
to low-carbon hours. Geographic: route to cleaner grids.'' Key result:
``Google achieved 40\% carbon reduction by shifting workloads globally.
Zero model changes.''

% -- WARN: What students will get wrong on THIS topic
Common error: students think carbon-aware scheduling requires sacrificing
latency. Correct framing: batch training jobs are delay-tolerant --- shifting
by 6--12 hours costs nothing but saves 30--50\% carbon.

% -- FLEX: [CORE] + contingency
[CORE] Carbon-aware scheduling is the most actionable technique in the chapter.
IF SHORT: State Google's 40\% result and the two strategies.
}

\footnotesize
\begin{columns}[T]
  \begin{column}{0.48\textwidth}
    \textbf{Grid carbon intensity varies 100$\times$:}

    \vspace{0.05cm}
    \scriptsize
    \renewcommand{\arraystretch}{1.0}
    \begin{tabular}{@{}lrl@{}}
      \toprule
      \textbf{Source} & \textbf{g/kWh} & \textbf{Regions} \\
      \midrule
      Coal   & 820--1,200 & Poland, W.\ Virginia \\
      Gas    & 350--500   & Texas \\
      Solar  & 20--50     & California \\
      Hydro  & 10--30     & \textcolor{datastroke}{\textbf{Quebec}} \\
      Nuclear & 5--20     & \textcolor{datastroke}{\textbf{France}} \\
      \bottomrule
    \end{tabular}
  \end{column}
  \begin{column}{0.48\textwidth}
    \textbf{Two strategies:}
    \begin{enumerate}\setlength\itemsep{0pt}
      \item \textbf{Temporal}: Delay to low-carbon hours
      \item \textbf{Geographic}: Route to cleaner grids
    \end{enumerate}

    \vspace{0.1cm}
    \begin{mlsyscard}{datastroke}
    {\scriptsize Google achieved \textbf{40\% carbon reduction} by shifting workloads globally. \textbf{Zero model changes.}}
    \end{mlsyscard}
  \end{column}
\end{columns}

\end{frame}

% --- ACTIVE LEARNING 3: Discussion ---
\begin{frame}{Discussion: Jevons Paradox}
\note{
% -- LINK: What prior concept connects to this slide
Optimization techniques reduce per-query cost. Jevons Paradox asks: does
cheaper AI actually reduce total consumption?

% -- NARRATE: What to SAY while showing this slide
[3 min] Point to the diagram: ``GPT-3 went from \$0.06 to \$0.002 per 1K
tokens --- 30x cheaper. This triggered 100x more queries. Net result: 3.3x
MORE total energy.'' Read the discussion prompt. Give 90 seconds for pairs.
Cold-call 2--3 pairs.

% -- ENGAGE: Specific question, prediction, or task for THIS slide
``If making AI cheaper increases total consumption, how do we actually reduce
emissions?'' [Expected answers: carbon budgets, caps, regulation. The point:
efficiency alone is insufficient without absolute limits.]

% -- FLEX: [CORE] + contingency
[CORE] Third active learning moment. Jevons Paradox is the chapter's most
counterintuitive insight.
IF SHORT: Reduce to 60 seconds discussion.
}

% --- Layout: FULL-WIDTH IMAGE ---
\centering
\safeimg[width=0.90\textwidth,height=4.0cm,keepaspectratio]{jevons-paradox.pdf}

\vspace{0.1cm}
{\large\bfseries Turn and Talk \textcolor{midgray}{(90 seconds)}}

\vspace{0.2cm}
{\normalsize\alert{If making AI cheaper increases total consumption, how do we actually reduce emissions?}}

\end{frame}

% =============================================================================
\section{Google's 4 Ms}
% =============================================================================

\begin{frame}{Google's ``4 Ms'' Framework}
\note{
% -- LINK: What prior concept connects to this slide
Individual techniques were covered separately. Google's 4 Ms framework shows
how they multiply together.

% -- NARRATE: What to SAY while showing this slide
[3 min] Walk through the diagram: ``Model: efficient architectures. Machine:
TPUs over GPUs. Mechanization: compiler and runtime optimization. Map: carbon-
aware geographic placement.'' Key result: ``Combined: 83x energy reduction,
747x carbon reduction on Transformer. These multiply, not add.''

% -- ENGAGE: Specific question, prediction, or task for THIS slide
Ask: ``Which M has the largest individual impact?'' [Expected: Map (geography)
--- 40x alone exceeds any single algorithmic technique.]

% -- FLEX: [CORE] + contingency
[CORE] The 4 Ms framework is the chapter's actionable summary.
IF SHORT: State the 747x combined result and name the four Ms.
}

% --- Layout: FULL-WIDTH IMAGE ---
\centering
\safeimg[width=0.92\textwidth,height=4.8cm,keepaspectratio]{four-ms-framework.pdf}

\vspace{0.1cm}
\small
\textbf{Combined: 83$\times$ energy, 747$\times$ \COtwo{}} (Transformer on TPU in clean region)
\mlsyscite{Patterson et al., ``Carbon Emissions and Large Neural Networks,'' 2022}

\end{frame}

\begin{frame}{Quantified Optimization Stack}
\note{
% -- LINK: What prior concept connects to this slide
The 4 Ms provided the framework. This slide quantifies the percentage savings
from each technique in the optimization stack --- replacing generic best practices
with measured numbers.

% -- NARRATE: What to SAY while showing this slide
[2 min] Walk through the table column by column: ``Region selection: 40x carbon
reduction, zero model changes. Quantization: 4--16x energy reduction. Pruning:
2--5x with 90\% weight removal. Distillation: 40\% fewer parameters. Transfer
learning: 20--100x vs.\ training from scratch.'' Key insight: ``These multiply ---
the 747x from Google's 4 Ms comes from applying all layers together.''

% -- FLEX: [CORE] + contingency
[CORE] Replaces generic guidelines with quantified optimization stack.
IF SHORT: Point to the top 3 rows and the multiplicative insight.
}

\footnotesize
\renewcommand{\arraystretch}{1.15}
\begin{tabular}{@{}llrl@{}}
  \toprule
  \textbf{Technique} & \textbf{CMOS Term} & \textbf{Savings} & \textbf{Effort} \\
  \midrule
  \textbf{Region selection}   & --- (grid carbon)  & \textbf{40$\times$} \COtwo{}  & Zero (config change) \\
  \textbf{Quantization} (INT8) & $C$, $V$          & 4--16$\times$ energy & Low (post-training) \\
  \textbf{Pruning} (90\%)      & $C$ (fewer gates)  & 2--5$\times$ energy  & Medium (retrain) \\
  \textbf{Distillation}        & All (smaller model) & 40\% fewer params   & High (teacher needed) \\
  \textbf{Transfer learning}   & All (skip training) & 20--100$\times$     & Low (fine-tune only) \\
  \textbf{Hardware selection}   & $\alpha$, $V$      & 5--13$\times$       & Medium (procurement) \\
  \textbf{Extend HW life} (3$\to$5yr) & Embodied   & 40\% embodied       & Organizational \\
  \bottomrule
\end{tabular}

\vspace{0.1cm}
\begin{mlsyscard}{datastroke}
{\footnotesize \textbf{These multiply, not add}: Region (40$\times$) $\times$ Hardware (5$\times$) $\times$ Quantization (4$\times$) = \textbf{800$\times$} carbon reduction before touching the model architecture.}
\end{mlsyscard}

\end{frame}

\begin{frame}{Regulatory Landscape}
\note{
% -- LINK: What prior concept connects to this slide
Engineering guidelines are voluntary. Regulation makes sustainability mandatory.

% -- NARRATE: What to SAY while showing this slide
[2 min] Point to the left: ``EU AI Act: models over 10 to the 25 FLOPs are
classified as systemic risk, requiring mandatory energy reporting. Fines: up
to 7\% of global revenue.'' Point to the right: ``CSRD mandates audited Scope
1, 2, 3 emissions for 50,000+ companies --- including GPU manufacturing carbon.''
Key point: ``Companies underreport Scope 3 by 50--70\%.''

% -- WARN: What students will get wrong on THIS topic
Common error: students assume sustainability reporting is optional or
aspirational. Correct framing: EU fines are cumulative and percentage-based ---
this is a hard compliance requirement.

% -- FLEX: [OPTIONAL] + contingency
[OPTIONAL] Regulatory context. IF SHORT: State the EU AI Act threshold and
the 7\% fine.
}

\footnotesize
\begin{columns}[T]
  \begin{column}{0.48\textwidth}
    \textbf{EU AI Act (2024):}
    \begin{itemize}\setlength\itemsep{0pt}
      \item Models $> 10^{25}$ FLOPs = ``systemic risk''
      \item Mandatory energy reporting (train + inference)
      \item Fines: up to \textbf{7\% of global revenue}
    \end{itemize}

    \vspace{0.1cm}
    \textbf{CSRD:}
    \begin{itemize}\setlength\itemsep{0pt}
      \item 50,000+ companies must disclose
      \item Audited Scope 1, 2, 3 emissions
      \item Includes GPU manufacturing carbon
    \end{itemize}
  \end{column}
  \begin{column}{0.48\textwidth}
    \textbf{GHG Protocol Scopes:}

    \vspace{0.1cm}
    \renewcommand{\arraystretch}{1.1}
    \begin{tabular}{@{}lr@{}}
      \toprule
      \textbf{Scope} & \textbf{Share} \\
      \midrule
      1 (Direct)     & 5--15\% \\
      2 (Grid power) & 60--75\% \\
      3 (Supply chain) & 15--25\% \\
      \bottomrule
    \end{tabular}

    \vspace{0.1cm}
    \begin{mlsyscard}{errorstroke}
    {\scriptsize Companies underreport Scope 3 by \textbf{50--70\%}.}
    \end{mlsyscard}
  \end{column}
\end{columns}

\end{frame}

% =============================================================================
\section{Wrap-Up}
% =============================================================================

\begin{frame}{Fallacies}
\note{
% -- NARRATE: What to SAY while showing this slide
[2 min] Read each fallacy with its quantitative rebuttal. The Jevons Paradox
fallacy is the most counterintuitive: 30x cheaper triggered 100x more queries
for 3.3x MORE total energy. The carbon offsets fallacy: 60--90\% of voluntary
credits fail to deliver claimed reductions.

% -- FLEX: [CORE] + contingency
[CORE] Fallacies correct dangerous misconceptions.
IF SHORT: Read fallacies 2 (Jevons) and 3 (offsets) only.
}

\small
\textbf{Fallacy:} \textit{Cloud computing automatically makes AI sustainable.}\\
{\footnotesize Geographic region dominates: same training run = 21$\times$ difference (US avg vs Quebec). Default cloud regions waste 20--50$\times$ carbon.}

\vspace{0.1cm}
\textbf{Fallacy:} \textit{Efficiency improvements reduce total environmental impact.}\\
{\footnotesize Jevons Paradox: GPT-3 became 30$\times$ cheaper per query, triggering 100$\times$ more queries. Net: 3.3$\times$ MORE total energy.}

\vspace{0.1cm}
\textbf{Fallacy:} \textit{Carbon offsets solve the sustainability problem.}\\
{\footnotesize 60--90\% of voluntary credits fail to deliver claimed reductions. Direct migration to renewables is 2--3$\times$ cheaper and more effective.}

\end{frame}

\begin{frame}{Pitfalls}
\note{
% -- NARRATE: What to SAY while showing this slide
[2 min] Read each pitfall. The first is surprising: on Quebec's clean grid,
embodied carbon is 98\% of total --- manufacturing dominates. Extending hardware
life from 3 to 5 years saves 40\%. The second: a model pruned 40\% for training
savings but requiring 2x inference compute increases total emissions at scale.

% -- FLEX: [CORE] + contingency
[CORE] Pitfalls highlight lifecycle thinking.
IF SHORT: Cover the first pitfall (embodied carbon) only.
}

\footnotesize
\textbf{Pitfall:} \textit{Focusing only on operational energy while ignoring embodied carbon.}\\
A single H100 embodies 164 kg \COtwo{} from manufacturing. On Quebec's clean grid, embodied carbon is \textbf{98\%} of total emissions. Extending hardware life from 3 to 5 years saves \textbf{40\%} --- larger than most algorithmic optimizations.

\vspace{0.15cm}
\textbf{Pitfall:} \textit{Optimizing individual components without lifecycle analysis.}\\
A model pruned 40\% to save training energy but requiring 2$\times$ inference compute increases total emissions if it serves $>$100M queries. Edge deployment saving 60\% datacenter energy but requiring 10,000 devices adds 1,500--2,000 kg embodied carbon (10$\times$ cloud training).

\end{frame}

% --- RETRIEVAL PRACTICE ---

% --- MUDDIEST POINT ---
\begin{frame}{Muddiest Point}
\note{
% -- NARRATE: What to SAY while showing this slide
[2 min] ``Write the single concept from today you found most confusing.
Anonymous, one sentence.'' Collect responses. Likely confusions: Jevons
Paradox mechanics, embodied vs.\ operational carbon, or PUE calculation.

% -- FLEX: [CORE] + contingency
[CORE] Closes the feedback loop. Do not skip.
}

\centering
\vspace{1.0cm}
{\Large\bfseries What was the \alert{muddiest point} today?}

\vspace{0.8cm}
{\normalsize Write down the concept you found \textbf{most confusing}.}

\vspace{0.5cm}
{\small\textcolor{midgray}{Anonymous. One sentence. Submit before you leave.}}

\end{frame}

\begin{frame}{What Were the Key Ideas?}
\note{
% -- NARRATE: What to SAY while showing this slide
[2 min] ``Close your notes. Write down the 4 most important concepts from today.
90 seconds, no peeking.'' Walk around the room. Do NOT show next slide yet.

% -- FLEX: [CORE] + contingency
[CORE] Retrieval practice is the highest-impact learning intervention.
}

\centering
\vspace{1.5cm}
{\Large\bfseries Close your notes.}

\vspace{0.8cm}
{\large Write down the \textbf{4 most important concepts} from today.}

\vspace{0.8cm}
{\normalsize\textcolor{midgray}{90 seconds --- no peeking.}}

\end{frame}

\begin{frame}{Key Takeaways}
\note{
% -- NARRATE: What to SAY while showing this slide
[2 min] Reveal after retrieval. Walk through each bullet, anchoring on numbers:
350,000x compute growth, 40x geographic leverage, 1,000x hardware range, Jevons
Paradox 3.3x rebound, 747x from Google's 4 Ms combined.

% -- FLEX: [CORE] + contingency
[CORE] Synthesis slide. IF SHORT: Read bullets 1, 3, and 5.
}

\scriptsize
\begin{itemize}\setlength\itemsep{0pt}
  \item \textbf{Sustainability Paradox}: AI compute grew 350,000$\times$ (2012--2019); hardware efficiency only 2--5\%/year. The gap is physical.
  \item \textbf{Lifecycle Carbon}: Training (60--80\%), inference (15--25\%), manufacturing (5--15\%). Inference exceeds training within days.
  \item \textbf{Geography dominates}: Site selection creates 40$\times$ emission differences --- larger than any algorithmic optimization.
  \item \textbf{CMOS physics}: $P = \alpha C V^2 f + V I_{leak}$. Quantization and pruning reduce $C$ and $V$ directly.
  \item \textbf{Jevons Paradox}: Efficiency alone increases total consumption. Requires optimization + carbon budgets.
  \item \textbf{Google's 4 Ms}: Model + Machine + Mechanization + Map = 747$\times$ \COtwo{} reduction (multiplicative).
  \item \textbf{Regulation}: EU AI Act mandates energy reporting for frontier models. Compliance requirement, not optional.
\end{itemize}

\end{frame}

\begin{frame}{References}
\note{
% -- NARRATE: What to SAY while showing this slide
[1 min] Highlight Patterson+21 and Patterson+22 as the foundational carbon
accounting papers. Strubell+19 started the conversation. Students can photograph.

% -- FLEX: [OPTIONAL] + contingency
[OPTIONAL] Do not read every entry aloud.
}

\small
\mlsysref{Patterson+21}{Patterson et al. ``Carbon Emissions and Large Neural Networks.'' 2021.}
\mlsysref{Patterson+22}{Patterson et al. ``The Carbon Footprint of Machine Learning Training.'' 2022.}
\mlsysref{Strubell+19}{Strubell et al. ``Energy and Policy Considerations for NLP.'' ACL 2019.}
\mlsysref{Schwartz+20}{Schwartz et al. ``Green AI.'' CACM 2020.}
\mlsysref{Hinton+15}{Hinton et al. ``Distilling the Knowledge in a Neural Network.'' 2015.}
\mlsysref{Sanh+19}{Sanh et al. ``DistilBERT.'' NeurIPS Workshop 2019.}

\end{frame}

\begin{frame}{Next Lecture: Responsible AI}
\note{
% -- NARRATE: What to SAY while showing this slide
[1 min] ``We have ensured the fleet is physically viable. The final question:
whom does it serve?'' Point to the three cards: fairness, governance, ethics.
``Next lecture completes the transition from how to build to whom it serves.''

% -- FLEX: [CORE] + contingency
[CORE] Forward hook motivates the final content chapter.
}

\footnotesize
\begin{columns}[c]
  \begin{column}{0.30\textwidth}
    \centering
    \begin{mlsyscard}{computestroke}
    {\large\bfseries Fairness}\\[0.1cm]
    {\scriptsize Bias detection\\Equitable outcomes}
    \end{mlsyscard}
  \end{column}
  \begin{column}{0.30\textwidth}
    \centering
    \begin{mlsyscard}{routingstroke}
    {\large\bfseries Governance}\\[0.1cm]
    {\scriptsize Accountability\\Transparency}
    \end{mlsyscard}
  \end{column}
  \begin{column}{0.30\textwidth}
    \centering
    \begin{mlsyscard}{errorstroke}
    {\large\bfseries Ethics}\\[0.1cm]
    {\scriptsize Value alignment\\Societal impact}
    \end{mlsyscard}
  \end{column}
\end{columns}

\vspace{0.3cm}
\centering
{\normalsize A system that is technically sound can still cause social harm.}\\[0.1cm]
{\small\textbf{From \emph{how to build} the machine to \emph{whom it serves}.}}

\end{frame}



\appendix

\begin{frame}{Backup: Carbon Calculation for Cloud Providers}
\note{
% -- NARRATE: What to SAY while showing this slide
Backup. Per-provider carbon intensity table showing 19x difference by region
choice within the same cloud provider.

% -- FLEX: [OPTIONAL] + contingency
[OPTIONAL] Only show if students ask about specific cloud regions.
}

\footnotesize
\renewcommand{\arraystretch}{1.15}
\begin{tabular}{@{}llrr@{}}
  \toprule
  \textbf{Provider} & \textbf{Region} & \textbf{g CO$_2$/kWh} & \textbf{Relative} \\
  \midrule
  AWS & us-east-1 (Virginia)  & 380 & 1.0$\times$ \\
  AWS & ca-central-1 (Quebec) & 20  & 0.05$\times$ \\
  GCP & us-central1 (Iowa)    & 430 & 1.13$\times$ \\
  GCP & europe-north1 (Finland) & 130 & 0.34$\times$ \\
  Azure & Sweden Central       & 30  & 0.08$\times$ \\
  \bottomrule
\end{tabular}

\vspace{0.1cm}
{\footnotesize Same model, same provider: 19$\times$ carbon difference by region choice alone.}

\end{frame}

\begin{frame}{Backup: Jevons Paradox Extended}
\note{
% -- NARRATE: What to SAY while showing this slide
Backup. Extended Jevons Paradox data: 60x cheaper per query but 5,000x more
queries yields 5x more total energy by 2025 projections.

% -- FLEX: [OPTIONAL] + contingency
[OPTIONAL] Only show if the Jevons discussion generated strong engagement.
}

\footnotesize
\textbf{Jevons Paradox in AI:} Efficiency leads to lower cost leads to higher demand.

\vspace{0.15cm}
\renewcommand{\arraystretch}{1.15}
\begin{tabular}{@{}lrrr@{}}
  \toprule
  \textbf{Year} & \textbf{Cost/1K tokens} & \textbf{Daily queries} & \textbf{Daily energy} \\
  \midrule
  2023 (GPT-3.5) & \$0.06  & 10M    & 10,000 kWh \\
  2024 (GPT-4o)  & \$0.005 & 500M   & 25,000 kWh \\
  2025 (proj.)   & \$0.001 & 5B     & 50,000 kWh \\
  \bottomrule
\end{tabular}

\vspace{0.1cm}
{\footnotesize 60$\times$ cheaper per query, but 5,000$\times$ more queries. Net: 5$\times$ more energy.}

\end{frame}


\end{document}
